\section{Introduction}

Shape control is an appealing route for studying and designing quantum-dot
spectra. Quantum dots and quantum billiards provide a natural setting for this
question because confinement, symmetry, and boundary geometry can all affect
low-energy states~\cite{reimann2002quantumdots,kouwenhoven2001fewelectron,berry1981sinai,stockmann1990quantumchaos}.
Tight-binding billiards and lattice quantum-dot models make these effects
computationally explicit while retaining finite-size and boundary-realization
features that are absent from ideal continuum sketches~\cite{cuevas1996chaoticbilliards,ulcakar2022tightbindingbilliards}.
These ingredients make a square-lattice tight-binding model of superelliptic
quantum dots a useful controlled benchmark: the Hamiltonian is simple, the
geometry family is interpretable, and spectra and conductance can be computed
with standard sparse-eigensolver and quantum-transport tools~\cite{groth2014kwant,datta2005quantumtransport}.
At the same time, this simplicity creates a demanding test of any proposed shape
objective. If an apparent signal is already explained by size, area, aspect
ratio, symmetry, degeneracy, boundary pixelation, reference-solver uncertainty,
or lead coupling, then it should not be presented as inverse design or as a new
physical effect.

This article is a methodology-oriented limitations paper. It does not enlarge
the parameter space, repair unsuccessful directions, or promote a new objective.
Instead, it proposes and applies a baseline-first audit hierarchy to a fixed set
of existing calculations in one finite superellipse tight-binding benchmark.
The hierarchy asks whether an apparent shape or ML-assisted signal survives
successively simpler explanations: confinement scale, area and perimeter,
aspect ratio, isotropy and degeneracy, lattice boundary realization,
reference-solver uncertainty, and contact geometry.

This kind of audit is needed because modern shape-controlled and
ML-assisted workflows can make weak observables look successful. A surrogate
may interpolate a dataset well while still learning size, aspect-ratio, or
symmetry structure; inverse-design candidates may then inherit the same
baseline rather than reveal a new physical effect. Classical surrogate-modeling
practice therefore requires comparison against simple, structured alternatives,
and spatially or hierarchically structured data require validation splits that
match the intended use~\cite{bishop2006prml,hastie2009elements,forrester2008surrogate,queipo2005surrogate,roberts2017crossvalidation}.

Shape-controlled quantum dots are especially vulnerable to this failure mode.
Low-energy levels of confined systems are strongly governed by dimensional
scaling; near the bottom of a square-lattice tight-binding band, the dispersion
has the expected quadratic form. In addition, size, area, aspect ratio,
symmetry, finite-lattice boundary realization, and transport contacts can all
mimic or absorb shape dependence. A baseline-first hierarchy is therefore a
necessary diagnostic before claiming a shape effect, a positive ML result, or a
device-design rule.

The main message is conservative and protocol-specific. In the tested
square-lattice tight-binding model of finite superelliptic quantum dots,
several plausible shape, ML, continuum-reference, and transport objectives do
not support robust positive claims after baseline checks. This does not imply
that geometry-based quantum design, machine learning, finite-size analysis, or
transport studies are useless in general. The useful output is the documented
audit hierarchy and the associated failure modes in this controlled benchmark.

\begin{figure*}[t]
    \centering
    \includegraphics[width=0.95\textwidth]{model_and_baseline_schematic.pdf}
    \caption{
    Schematic of the closed finite-lattice superellipse dot, the corresponding
    open-dot transport geometry with separate left and right leads, and the
    baseline ladder used before interpreting shape-dependent observables. The
    ladder is conceptual: it describes the order of checks, not a new numerical
    calculation.
    }
    \label{fig:model-baseline-schematic}
\end{figure*}
